% Part:sets-functions-relations
% Chapter: size-of-sets
% Section: non-enumerability-alt

\documentclass[../../../include/open-logic-section]{subfiles}

\begin{document}

\olfileid[ar]{sfr}{siz}{nen-alt}

\olsection{مجموعات \printtoken{S}{nonenumerable}}

\begin{editorial}
نبين هنا أن $\Bin^\omega$ و$\Pow{\Nat}$ غير قابلتين للتعداد، ونعتمد تعريفات \olref[enm-alt]{sec}. فالمشترط دالة تقابلية مع~$\Nat$، لا دالة شاملة من~$\PosInt$.
\end{editorial}

\begin{explain}
مجموعة الأعداد الطبيعية~$\Nat$ لا نهائية، وكونها !!{enumerable} بين بنفسه. وإنما الأمر اللافت وجود مجموعات \emph{!!{nonenumerable}}، أي ليست !!{enumerable} (انظر \olref[sfr][siz][enm-alt]{defn:enumerable}).

وقد يستغرب هذا؛ فإن كون $A$ !!{nonenumerable} يقتضي \emph{انتفاء} كل !!{bijection} $f
\colon \Nat \to A$. بل لا توجد دالة ترسل !!{element}s~$\Nat$، على كثرتها التي لا تنتهي، إلى~$A$ وتستوفي~$A$. فإذا كانت $A$ !!{nonenumerable}، ففي~$A$ !!{element}s «أكثر» من الأعداد الطبيعية.

فسبيل إثبات أن المجموعة !!{nonenumerable} أن ننفي وجود !!{bijection} على الوجه المطلوب. ويكفي لذلك أن نبرهن أن كل محاولة لتعداد !!{element}s~$A$ تفوت !!{element} واحدًا على الأقل؛ فتكون كل دالة $f\colon \Nat \to A$ غير !!{surjective}. وطريقة كانتور \emph{القطرية} سبيل عام إلى هذا النفي: تعطى قائمة من !!{element}s~$A$، ولتكن $x_1$، $x_2$، \dots، فننشئ !!{element} من~$A$ يمتنع، بحكم إنشائه، أن يكون في القائمة.

وليتضح الأمر بالمثال، نبدأ بالمجموعة~$\Bin^\omega$ التي تضم جميع السلاسل اللانهائية من~$0$ و~$1$. (فالرمز «$\Bin$» يدل على الثنائية، ويكفينا أن نعده اسمًا للمجموعة ذات العنصرين $\{0,1\}$.)\oliflabeldef{sfr:card-arithmetic:card-opps:sec}{\footnote{والتحقيق أن $\Bin^\omega$ مجموعة جميع متتاليات $\omega$ المؤلفة من~$0$ و~$1$، أي المجموعة $\funfromto{\omega}{\{0,1\}}$. وإنما يتبين معنى هذا في \olref[sfr][card-arithmetic][card-opps]{sec}.}{} أما هذه العبارة التي فيها شيء من التجوز فلا يترتب عليها لبس فيما نحن فيه.}
\end{explain}

\begin{thm}
\ollabel{thm:nonenum-bin-omega}
$\Bin^\omega$~!!{nonenumerable}.
\end{thm}

\begin{proof}
خذ تعدادًا كيف اتفق لمجموعة جزئية من~$\Bin^\omega$. فتكون بين أيدينا قائمة $s_{0}$، $s_{1}$، $s_{2}$، \dots{}، وكل~$s_n$ سلسلة لا نهائية من~$0$ و~$1$. وللرمز $s_n(m)$ مؤشران: $m$ يحدد رتبة الرقم، و$n$ يحدد السلسلة. فنرسم القائمة مصفوفة يقع فيها $s_n(m)$ في الصف ذي الرتبة~$n$ والعمود ذي الرتبة~$m$:

\[
\begin{array}{c|c|c|c|c|c}
& 0 & 1 & 2 & 3 & \dots \\\hline
0 & \mathbf{s_{0}(0)} & s_{0}(1) & s_{0}(2) & s_0(3) & \dots \\\hline
1 & s_{1}(0)& \mathbf{s_{1}(1)} & s_1(2) & s_1(3) & \dots \\\hline
2 & s_{2}(0)& s_{2}(1) & \mathbf{s_2(2)} & s_2(3) & \dots \\\hline
3 & s_{3}(0)& s_{3}(1) & s_3(2) & \mathbf{s_3(3)} & \dots \\\hline
\vdots & \vdots & \vdots & \vdots & \vdots & \mathbf{\ddots}
\end{array}
\]
وننشئ سلسلة لا نهائية~$d$ من~$0$ و~$1$ لا ترد في هذه القائمة. وطريق إنشائها تعيين خاناتها كلها، أي تعيين $d(n)$ لكل $n \in
\Nat$. فنقرأ قطر المصفوفة---ومن هنا اسم الطريقة القطرية---ونستبدل كل~$1$ بـ~$0$، وكل~$0$ بـ~$1$. وضبط ذلك أن نجعل $d(n)$ مساويًا لـ$0$ أو~$1$ تبعًا لكون !!{element} ذي الرتبة~$n$ على القطر، أعني $s_n(n)$، مساويًا لـ$1$ أو~$0$:

\[
d(n) =
\begin{cases}
1 & \text{إذا كان $s_{n}(n) = 0$}\\
0 & \text{إذا كان $s_{n}(n) = 1$}
\end{cases}
\]
فيتحقق $d \in \Bin^\omega$؛ إذ هي سلسلة لا نهائية من~$0$ و~$1$. وقد عينّا~$d$ بحيث $d(n) \neq s_n(n)$ لكل $n \in
\Nat$؛ فهي، أي~$d$، تخالف~$s_n$ في الخانة ذات الرتبة~$n$. إذن $d \neq s_n$ لكل $n\in \Nat$، فلا تكون~$d$ في القائمة $s_0$، $s_1$، $s_2$، \dots

فمهما أعطينا تعدادًا لمجموعة جزئية من $\Bin^\omega$، وجدنا !!{element} من~$\Bin^\omega$ خارجًا عنه. فلا تعداد يستوفي المجموعة~$\Bin^\omega$؛ أي إن~$\Bin^\omega$ !!{nonenumerable}.
\end{proof}

\begin{editorial}
\textbf{تنبيه تصحيحي على الأصل (OLSIZ-008).}
عكس الأصل وصفي المؤشرين؛ فالجدول يضع $s_n(m)$ في الصف $n$
والعمود $m$، ولذلك هي الخانة $m$ من السلسلة $n$.

\textbf{تنبيه تصحيحي على الأصل (OLSIZ-009).}
كرر الأصل التحويل من $1$ إلى $0$ وأسقط التحويل المقابل؛
فأُثبت التحويلان $1\mapsto0$ و$0\mapsto1$ اللذان تعطيهما صيغة الحالات.
\end{editorial}

\begin{explain}
واسم «الاستدلال القطري» مأخوذ من استعمال قطر المصفوفة في تعريف~$d$. غير أن المصفوفة ليست من لوازمه؛ فقد نثبت أن مجموعة !!{nonenumerable} بمعنى مماثل، من غير مصفوفة مرسومة ولا قطر ظاهر. وبيان ذلك في النتيجة الآتية.
\end{explain}

\begin{thm}
\ollabel{thm:nonenum-pownat}
$\Pow{\Nat}$ ليست !!{enumerable}.
\end{thm}

\begin{proof}
نسلك المسلك نفسه، فنبين أن أي قائمة من المجموعات الجزئية من~$\Nat$ تفوتها مجموعة جزئية من~$\Nat$. ولتكن $N_0, N_1, N_2, \ldots$ قائمة مفروضة من المجموعات الجزئية من~$\Nat$. ونعين مجموعة~$D$ بهذه الصفة: $n \in D$ إذا وفقط إذا كان $n \notin
N_{n}$، أي
\[
D = \Setabs{n \in \Nat}{n \notin N_n}
\]
فإن $D\subseteq \Nat$ ظاهر. وأما~$D$ فلا ترد في القائمة؛ إذ $n \in D$ يكافئ $n\notin N_n$ بمقتضى التعريف، ومنه $D \neq N_n$ لكل $n \in \Nat$.
\end{proof}

\begin{explain}
لم نذكر قطرًا في البرهان، ولكن يظهر معناه إذا رسمنا المجموعات $N_0$، $N_1$، \dots{} في مصفوفة. فنجعل $N_n$ في الصف ذي الرتبة~$n$، ونكتب~$m$ في العمود ذي الرتبة~$m$ إذا وفقط إذا كان $m \in N_n$. ولو كانت المجموعات الأربع الأولى $\{0,1,2,\dots\}$، $\{1, 3, 5, \dots\}$، $\{0,1,4\}$، و$\{2,3,4,\dots\}$، لابتدأت المصفوفة هكذا:
\[
\begin{array}{r@{}rrrrrrr}
  N_0 = \{ & \mathbf{0}, & 1, & 2, & & & & \dots\}\\
  N_1 = \{ &  & \mathbf{1}, &  & 3, &  & 5, & \dots\}\\
  N_2 = \{ & 0, & 1, &  &  & 4\phantom{,} &  & \}\\
  N_3 = \{ &  &  & 2, & \mathbf{3}, & 4, & & \dots\}\\
  &\vdots & & & & & & \ddots\phantom{\}}
  \end{array}
\]
فنحصل على~$D$ بالنزول على القطر، جاعلين $n \in D$ إذا وفقط إذا كان~$n$ \emph{غائبًا} عن موضعه القطري. ففي هذا المثال نستبعد~$0$ و~$1$، ونثبت~$2$، ونستبعد~$3$، ثم نمضي على ذلك.
\end{explain}

\begin{prob}
أقم حجة قطرية صريحة تبين أن مجموعة جميع الدوال $f \colon \Nat \to
\Nat$ !!{nonenumerable}. أي بين أن كل قائمة $f_1$، $f_2$، \dots{} من الدوال التي نوع كل منها $f_i\colon \Nat \to \Nat$ تفوتها دالة $g \colon \Nat \to \Nat$ من النوع نفسه.
\end{prob}

\end{document}
